\documentclass[11pt,a4paper]{article}

\usepackage[margin=25mm]{geometry}
\usepackage[T1]{fontenc}
\usepackage{lmodern}
\usepackage{microtype}
\usepackage{amsmath,amssymb}
\usepackage{graphicx}
\usepackage{booktabs}
\usepackage{array}
\usepackage{float}
\usepackage{subcaption}
\usepackage{siunitx}
\usepackage{xcolor}
\usepackage{listings}
\usepackage[hidelinks]{hyperref}

\sisetup{per-mode=symbol,detect-all}
\graphicspath{{figures/}}
\setlength{\emergencystretch}{2em}
\definecolor{codeblue}{RGB}{31,78,121}
\definecolor{codegray}{RGB}{245,245,245}
\lstset{
  language=Python,
  basicstyle=\ttfamily\small,
  keywordstyle=\color{codeblue}\bfseries,
  backgroundcolor=\color{codegray},
  frame=single,
  breaklines=true,
  showstringspaces=false
}

% The report compiles before the notebooks have been run.  A labelled box is
% shown for a missing plot and is replaced automatically after figure export.
\newcommand{\simfigure}[2][0.94]{%
  \IfFileExists{#2}{%
    \includegraphics[width=#1\linewidth]{#2}%
  }{%
    \fbox{\parbox[c][48mm][c]{0.88\linewidth}{\centering
      Run the notebooks to generate\\[2mm]
      \texttt{\detokenize{#2}}}}%
  }%
}

\title{Primary Exploration of Python Packages Optiland and xrt for Ray Optics.}
\author{}
\date{\today}

\begin{document}
\maketitle


\section{Division of software responsibilities.}
Optiland is used for sequential visible-light surfaces, glass dispersion,
reflection, and image-plane tracing.  xrt/raycing is used for the geometric
source--mirror--screen beamline, polarization weights, finite apertures, and
X-ray material reflectivity.  In the lens notebook, xrt's documented external
refractive-index interface is used as a material bridge.  Native xrt refractive
elements are aimed primarily at compound refractive X-ray lenses, so ordinary
visible spherical lenses are traced by Optiland rather than being forced into
an inappropriate X-ray model.

\section{Plane-mirror simulation}

\subsection{Default configuration}

The supplied example is an opaque gold mirror at a wavelength of
\SI{0.15406}{\nano\metre}.  The equivalent photon energy is
\begin{equation}
 E_\gamma=\frac{hc}{\lambda}
 =\frac{\SI{1239.841984}{\electronvolt\nano\metre}}
 {\SI{0.15406}{\nano\metre}}
 \simeq \SI{8.048}{\kilo\electronvolt}.
\end{equation}
This default lies in the intended range of the xrt scattering-factor tables.
For a visible or infrared mirror, the notebook requires the constant mode and
the measured coating values of $R$ and $A$.

\begin{table}[H]
\centering
\caption{Principal default parameters in \texttt{MirrorConfig}.}
\begin{tabular}{@{}lll@{}}
\toprule
Parameter & Default & Role \\
\midrule
Wavelength & \SI{0.15406}{\nano\metre} & photon energy/material response \\
Grazing angle & \SI{0.50}{\degree} & angle from the xrt mirror surface \\
Source--mirror distance & \SI{1000}{\milli\metre} & incoming propagation \\
Mirror--screen distance & \SI{800}{\milli\metre} & outgoing propagation \\
Source size $(x,z)$ & $\SI{0.40}{\milli\metre}\times\SI{0.20}{\milli\metre}$ & footprint size \\
Divergence $(x,z)$ & $(0.10,0.05)\,\si{\milli\radian}$ & angular broadening \\
Mirror material & Au, \SI{19.32}{\gram\per\cubic\centi\metre} & xrt optical constants \\
Rays & $20\,000$ & Monte-Carlo sample size \\
\bottomrule
\end{tabular}
\end{table}

\subsection{Geometry and reflection invariant}

For incident unit direction $\hat{\mathbf d}$ and unit surface normal
$\hat{\mathbf n}$, the reflected direction is
\begin{equation}
 \hat{\mathbf d}_{r}=\hat{\mathbf d}
 -2(\hat{\mathbf d}\cdot\hat{\mathbf n})\hat{\mathbf n}.
 \label{eq:reflection}
\end{equation}
Thus a mirror rotation by $\Delta\theta$ rotates the specular beam by
$2\Delta\theta$.  In the default xrt layout the expected vertical deflection
is therefore \SI{1.00}{\degree}.  The screen centre is placed on this direction,
so its elevation relative to the mirror is
\begin{equation}
 z_s=L\sin(2\theta)
 =\SI{800}{\milli\metre}\sin(\SI{1.00}{\degree})
 \simeq\SI{13.96}{\milli\metre}.
\end{equation}
The notebook asserts the unit length of the reflected direction and checks the
mean Monte-Carlo deflection against $2\theta$.

The Optiland illustration uses the same vector law in a visually convenient
folded geometry: a \SI{45}{\degree} plane mirror turns a beam initially along
$+z$ into $+y$.  Its detector is positioned normal to that reflected branch.

\begin{figure}[H]
\centering
\begin{subfigure}[t]{0.49\textwidth}
  \centering
  \simfigure{figures/plane_mirror_optiland.png}
  \caption{Optiland fold-mirror layout.}
\end{subfigure}\hfill
\begin{subfigure}[t]{0.49\textwidth}
  \centering
  \simfigure{figures/plane_mirror_layout.png}
  \caption{xrt Monte-Carlo scattering plane.}
\end{subfigure}
\caption{The two geometric views implement the same specular-reflection law
but use different coordinate conventions.}
\label{fig:mirror-layouts}
\end{figure}

\subsection{Reflectance, absorptance, and detector distribution}

xrt obtains the complex $s$- and $p$-polarized amplitudes from the selected
material and incidence angle.  The intensity reflectances are
\begin{equation}
 R_s=|r_s|^2,\qquad R_p=|r_p|^2.
\end{equation}
These weights are applied to the ray coherency matrix.  For an opaque thick
mirror, transmission is neglected and
\begin{equation}
 A=1-R,\qquad T=0,\qquad R+A+T=1.
 \label{eq:budget}
\end{equation}
In the user-defined constant mode, all three branches are retained in the
reported budget, with $T=1-R-A$.  xrt traces the reflected branch geometrically;
the table and bar chart account for the non-reflected power.

Figure~\ref{fig:mirror-results} should be interpreted as follows.  The
reflectance plot shows the rapid angle and polarization dependence expected of
an X-ray mirror.  The vertical dashed line marks the selected grazing angle.
The detector distribution convolves source size and divergence and is clipped
by the finite mirror.  Increasing either divergence broadens the corresponding
screen coordinate approximately linearly with propagation distance; decreasing
the mirror width or length lowers the accepted-ray fraction.

\begin{figure}[H]
\centering
\begin{subfigure}[t]{0.49\textwidth}
  \centering
  \simfigure{figures/plane_mirror_reflectance.png}
  \caption{Polarization-dependent material reflectance.}
\end{subfigure}\hfill
\begin{subfigure}[t]{0.49\textwidth}
  \centering
  \simfigure{figures/plane_mirror_screen.png}
  \caption{Reflected screen footprint.}
\end{subfigure}
\begin{subfigure}[t]{0.62\textwidth}
  \centering
  \simfigure{figures/plane_mirror_energy_budget.png}
  \caption{Reflected, absorbed, and transmitted fractions.}
\end{subfigure}
\caption{Plane-mirror outputs generated by \texttt{Plane Mirror.ipynb}.}
\label{fig:mirror-results}
\end{figure}

\section{Movable lens-group simulation}

\subsection{Prescription and exact model}

The default three-element group is summarized in Table~\ref{tab:lenses}.  The
front and rear radii, thickness, glass, position, wavelengths, field angles,
aperture, and fixed sensor location are all editable.  Removing either the
second or third tuple produces a two-lens model without changing the builder.

\begin{table}[H]
\centering
\caption{Default front-vertex prescription and paraxial reference values.}
\label{tab:lenses}
\begin{tabular}{@{}lrrrrl@{}}
\toprule
Lens & $z$ (mm) & $R_1$ (mm) & $R_2$ (mm) & $d$ (mm) & Glass \\
\midrule
L1 & 0  & 100 & $-100$ & 5 & N-BK7 \\
L2 & 45 & 150 & $-150$ & 4 & N-SF5 \\
L3 & 70 & 120 & $-120$ & 5 & N-BK7 \\
\bottomrule
\end{tabular}
\end{table}

For the first-order comparison, each element focal length is estimated with the
thick lensmaker equation in air,
\begin{equation}
 \frac{1}{f}=(n-1)\left[
 \frac{1}{R_1}-\frac{1}{R_2}
 +\frac{(n-1)d}{nR_1R_2}\right].
 \label{eq:lensmaker}
\end{equation}
The reference indices in the notebook give approximately
$f_1=\SI{97.58}{\milli\metre}$,
$f_2=\SI{112.09}{\milli\metre}$, and
$f_3=\SI{116.93}{\milli\metre}$.  Optiland does not use these constants for
the exact trace: it loads the named glass and evaluates its dispersion at each
of \SI{486.1}{\nano\metre}, \SI{587.6}{\nano\metre}, and
\SI{656.3}{\nano\metre}.

\begin{figure}[H]
\centering
\begin{subfigure}[t]{0.49\textwidth}
  \centering
  \simfigure{figures/lens_group_baseline.png}
  \caption{Baseline Optiland surfaces.}
\end{subfigure}\hfill
\begin{subfigure}[t]{0.49\textwidth}
  \centering
  \simfigure{figures/lens_group_moved.png}
  \caption{Lens 2 translated by \SI{10}{\milli\metre}.}
\end{subfigure}
\caption{Exact thick-lens layouts at the same fixed sensor plane.}
\label{fig:lens-layouts}
\end{figure}

\subsection{Why lens translation moves the focus}

In one meridional plane, free propagation by distance $d$ and refraction by an
ideal thin lens of focal length $f$ are represented by
\begin{equation}
 P(d)=\begin{pmatrix}1&d\\0&1\end{pmatrix},\qquad
 L(f)=\begin{pmatrix}1&0\\-1/f&1\end{pmatrix}.
\end{equation}
The group matrix contains alternating $P$ and $L$ factors.  Moving lens 2
changes both adjacent propagation distances, so the group power and the
principal-plane locations change even though the lens curvature does not.
For collimated input, a matrix
$M=\left(\begin{smallmatrix}A&B\\C&D\end{smallmatrix}\right)$ places the
paraxial focus a distance $-A/C$ after the output reference plane.

With the default reference indices, the paraxial calculation predicts a focus
near $z=\SI{79.88}{\milli\metre}$, close to the fixed sensor at
\SI{80}{\milli\metre}.  Translating L2 from $z=\SI{45}{\milli\metre}$ to
$z=\SI{55}{\milli\metre}$ moves the first-order focus to about
$z=\SI{83.96}{\milli\metre}$.  Therefore the moved configuration should have a
larger spot at the unchanged sensor.  Exact numerical values in the Optiland
table can differ because the spherical surfaces include thickness, aberration,
field angle, vignetting, and wavelength-dependent glass data.

\begin{figure}[H]
\centering
\simfigure[0.90]{figures/lens_group_comparison.png}
\caption{Independent paraxial trajectories before and after translating L2.}
\label{fig:lens-comparison}
\end{figure}

\subsection{Spot diagrams and position sensitivity}

The image-plane RMS spot radius is calculated from accepted rays as
\begin{equation}
 r_{\mathrm{RMS}}=
 \sqrt{\left\langle(x-\bar{x})^2+(y-\bar{y})^2\right\rangle}.
 \label{eq:rms}
\end{equation}
A smaller value at the fixed sensor means better convergence there, not
necessarily a globally optimized imaging system.  Differences between the
three wavelength panels show chromatic residuals.  At the off-axis field,
coma, astigmatism, and field curvature may also contribute.

The two-dimensional sweep moves L2 and L3 while keeping L1 and the sensor
fixed.  A dark valley indicates combinations that make the first-order sensor
height small.  The sweep is an initializer: promising positions must be copied
back to the prescription and verified with the exact Optiland trace.  Invalid
overlapping prescriptions are masked rather than silently traced.

\begin{figure}[H]
\centering
\begin{subfigure}[t]{0.59\textwidth}
  \centering
  \simfigure{figures/lens_group_spots.png}
  \caption{Exact multi-wavelength spot distributions.}
\end{subfigure}\hfill
\begin{subfigure}[t]{0.39\textwidth}
  \centering
  \simfigure{figures/lens_group_sweep.png}
  \caption{Paraxial lens-position sensitivity.}
\end{subfigure}
\caption{Fixed-sensor convergence metrics from \texttt{lens\_group.ipynb}.}
\label{fig:lens-results}
\end{figure}

\section{Validation, interpretation, and limitations}

The notebooks use several checks that are independent of a visually plausible
plot:
\begin{enumerate}
  \item Equation~\eqref{eq:reflection} predicts the $2\theta$ mirror
  deflection, which is compared with the mean xrt ray direction.
  \item Equation~\eqref{eq:budget} is asserted numerically, including a
  possible transmitted fraction in constant mode.
  \item The external visible refractive indices returned by xrt are checked
  against the requested constants.
  \item The paraxial lens calculation predicts the direction of focus shift;
  the exact Optiland RMS spot verifies its consequence at the fixed sensor.
  \item Lens overlap and sensor-order checks reject non-physical position
  combinations during the sweep.
\end{enumerate}

The simulations are geometrical-optics models.  They do not by themselves
include diffraction-limited point-spread functions, mirror roughness scatter,
surface figure error, mechanical tolerances, thermal deformation, ghost paths,
or a measured coating spectrum.  The xrt material tables are suitable for the
X-ray/EUV example, not for an arbitrary visible mirror.  Conversely, visible
glass catalogue data should not be interpreted as an X-ray compound refractive
lens model.  For an experiment, measured source phase space, coating data,
surface prescriptions, and detector response should replace the defaults.

\section{Recommended parameter studies}

For the mirror, sweep wavelength and grazing angle together, plot $R_s$ and
$R_p$, and record accepted flux and screen RMS width.  This separates a
material-response loss from aperture clipping.  For the lens group, first use
the fast position heat map to select a region, then evaluate the exact
multi-wavelength RMS spot and off-axis fields at each candidate.  A final design
should optimize radii, spacings, glass choices, and sensor position jointly
rather than interpreting one translated configuration as an optimum.

\begin{thebibliography}{9}
\bibitem{optiland}
Optiland documentation, ``Quickstart---Your First 5 Minutes'' and reflective
examples.  \url{https://optiland.readthedocs.io/}

\bibitem{xrt}
xrt documentation, ``Raycing backend'' and ``Using xrt as X-ray calculator.''
\url{https://xrt.readthedocs.io/}
\end{thebibliography}

\end{document}
